% formal/flow-graph-ch2021-comparison.tex -- CH2021 versus this project
%
% Status: source-routed comparison scaffold, not a proof.  Mined from the
% stale flow-graph-formal-axioms worktree on 2026-06-22.
%
% Purpose:
%   Record which parts of CH2021 are inspiration or terminology for the
%   flow-graph/tube work, and which parts are not imported as correctness
%   arguments.  This file is a guardrail for future formal and thesis writing.
%
% CH2021 source anchors checked:
%   papers/ch2021/s1_introduction_and_main_results.tex
%     def:cro, Type 1/2/3 definition, thm:combtosmooth, thm:smoothtocomb,
%     cor:computecehz.
%   papers/ch2021/s2_type_1_reeb_orbits.tex
%     def:linear_flow_graph, def:flow_graph_periodic_orbit,
%     def:sfgp, prop:orbitbijection, the footnote after prop:orbitbijection.
%   papers/ch2021/s6_smooth_combinatorial.tex
%     lem:combtosmooth and the proof of thm:smoothtocomb.
%
% Project anchors:
%   formal/flow-graph-real-algorithm.tex
%     the current idealized real-number flow-graph/tube theorem surface.
%   formal/flow-graph-capacity.tex
%     def:flow-graph-capacity-search-axioms,
%     thm:abstract-flow-graph-capacity,
%     rem:flow-graph-proof-obligations.
%   crates/symplectic/src/algorithms/flow_graph/README.md.

% input by formal/main.tex -- no \documentclass

\subsection{CH2021 Comparison}
\label{sec:flow-graph-ch2021-comparison}

CH2021 is a paper about combinatorial Reeb dynamics on symplectic
four-polytopes, smoothings, rotation numbers, and finite computations of
Reeb-orbit data.  This project uses some CH2021 language and geometric ideas,
but its flow-graph capacity theorem is not a specialization of the CH2021
capacity corollary.  The project route is the abstract route in
Theorem~\ref{thm:abstract-flow-graph-capacity}: HK2017 supplies a simple
capacity-achieving orbit, and the remaining work is to prove soundness and
completeness of the project's closed-word/tube resolver.

\begin{remark}[Non-transfer rule]\label{rem:ch2021-non-transfer-rule}
Do not prove a thesis-facing flow-graph capacity claim by saying that CH2021 has
a flow-graph algorithm and this project has a flow-graph implementation.  The
shared phrase ``flow graph'' is not a proof bridge.  A CH2021 statement can be
used only after checking that the objects, hypotheses, and output claim match
the project statement being made.
\end{remark}

\begin{unverified}
\begin{definition}[Comparison categories]\label{def:flow-graph-comparison-categories}
For this comparison:
\begin{itemize}
  \item \emph{Paper-side object} means the object as stated in CH2021 source
    files.
  \item \emph{Project-side object} means the object used by the current
    flow-graph formal scaffold or Rust work surface.
  \item \emph{Imported} means the statement can be used as a source-backed
    dependency after citation and notation translation.
  \item \emph{Inspired} means the statement motivates terminology or design but
    does not prove a project theorem.
  \item \emph{Not used} means future thesis text should avoid making the
    project depend on that statement unless a new explicit bridge is written.
\end{itemize}
\end{definition}
\end{unverified}

\begin{remark}[Comparison points]\label{rem:ch2021-project-comparison}
The differences that matter for proof reuse are:
\begin{description}
  \item[Main computational target.]
    CH2021 finds combinatorial Reeb orbits and actions, and for Type~1 orbits
    also defines rotation/CZ data.  Type~2 rotation and CZ data are not
    defined.  The paper
    uses these to test Viterbo-related questions and compute capacity through
    the CH2021 smooth-combinatorial route.  This project's flow-graph route
    targets full EHZ capacity only under the supported exact theorem scope
    stated explicitly in this project. Its correctness theorem is organized by
    Definition~\ref{def:flow-graph-capacity-search-axioms}.
  \item[Capacity proof route.]
    CH2021's capacity corollary uses the smooth-to-combinatorial theorem,
    rotation bound, and Type 1/Type 2 alternatives.  This project uses
    Theorem~\ref{thm:abstract-flow-graph-capacity}: HK2017 simple minimizers
    plus resolver soundness/completeness.  CH2021's capacity corollary is not a
    load-bearing input.
  \item[Flow-graph nodes.]
    In CH2021's 4d polytope flow graph, nodes are 2-faces and domains are those
    2-faces.  In this project, closed words are facet words; primitive tubes use
    triples of facets and have endpoint polygons on two-face charts.
  \item[Periodic-orbit condition.]
    In CH2021, a graph periodic orbit is a cycle plus a fixed point of the
    composed affine flow map.  In this project, a closed word is resolved by
    tube construction and a fixed-point problem on the starting two-face.  These
    have the same broad fixed-point shape, but the resolver semantics are
    project-specific.
  \item[Type 1/2/3 classification.]
    CH2021 classifies combinatorial Reeb orbits into Type 1, Type 2, and Type
    3.  Type 1 avoids the 1-skeleton; Type 2 meets it only at finitely many
    segment endpoints; Type 3 contains a bad 1-face.  The classification is not
    part of the project capacity theorem.  The project must not introduce a
    Type 2 obligation unless it explicitly switches to a CH2021-style capacity
    claim.
  \item[Finite search control.]
    CH2021 controls the capacity search with the rotation-derived bound
    \(\sum_F c_F\leq2\), which bounds the relevant segment count; an action
    cutoff alone does not bound arbitrarily long helical words.  The project
    search control is expressed through a finite
    supported search domain~$D_K$, lower-bound certificates, and concrete
    search proofs.
  \item[Numerical implementation.]
    CH2021 mentions computer implementation optimizations, but the paper
    statements are not this repository's API contract. The current Rust path
    distinguishes exact accepted output, typed rejection, retained words, and
    profiling/evidence surfaces. The earlier binary64 prototype is retired.
\end{description}
CH2021's statement that symplectic polytopes are generic is distinct from the
project's determinant-generic theorem.  In particular, density of the
paper's stronger condition excluding Type~2 orbits below an action bound is
conjectural there, not an imported open-dense result for this project.
\end{remark}

\begin{remark}[What CH2021 can still contribute]\label{rem:ch2021-useful-inputs}
The following CH2021 material remains useful for this project:
\begin{enumerate}
  \item the language of linear flow graphs, affine flow maps, composed action
    functions, and fixed points of composed maps;
  \item the geometric picture that transitions between 3-face Reeb segments can
    be organized by two-face data;
  \item examples and figures that help explain why a flow-graph picture is
    natural;
  \item warnings about lower-dimensional faces and degeneracies.
\end{enumerate}
Each use still needs a local bridge from CH2021 notation to the project's
facet-word/tube notation before it supports theorem-strength text.
\end{remark}

\begin{remark}[What this project must prove locally]\label{rem:ch2021-local-obligations}
The project must prove or explicitly assume the following locally:
\begin{enumerate}
  \item the supported closed-word/tube resolver is sound: every returned action
    is the action of a generalized closed Reeb orbit;
  \item the resolver is complete for the HK2017 simple capacity-achieving orbit
    on the supported input class;
  \item the finite supported search domain~$D_K$ contains the relevant simple
    capacity-achieving word;
  \item cached construction and action-cutoff search preserve the lower-bound
    certificates required by Definition~\ref{def:flow-graph-capacity-search-axioms};
  \item any numerical frontend used for accepted output agrees with the exact
    semantic resolver, or all output-relevant numerical predicates are resolved
    exactly. This is a general soundness requirement, not a current FG
    implementation claim.
\end{enumerate}
These are project obligations even when their geometric motivation came from
CH2021.
\end{remark}

\begin{remark}[Thesis-writing guardrail]\label{rem:ch2021-thesis-guardrail}
In thesis prose, it is safe to say that the algorithm is inspired by CH2021's
flow-graph description of Type 1 combinatorial Reeb orbits, but only as
geometric language and motivation, not as a Type 1 orbit model for this
project.  It is not safe to say that CH2021 proves this implementation
computes $c_{\mathrm{EHZ}}$ unless the thesis explicitly states and proves a
bridge from the CH2021 hypotheses and orbit classification to this project's
supported search domain and resolver.
\end{remark}
